\hspace{3mm}

\centerline{\bf \Huge  Удивительная задача}

\begin{flushright}
    \scriptsize{Прочитал неправильно условие $-$ ошибка\\ Не решил задачу с новым условием $-$ фатальная ошибка}
\end{flushright}

\textbf{\large Оригинальная постановка задачи.}

\problem{1}
$\bigl[\textbf{Олимпиада Курчатов, 2024, 11.3}\bigr]$ Аня и Боря играют в игру. Они по очереди (начинает Аня) выписывают по одной цифре, пока не получится шестизначное число. При этом первая выписанная цифра ненулевая и все выписанные цифры различны. Аня выигрывает, если полученное шестизначное число делится хотя бы на одно из чисел: 2, 3 или 5. Если этого не случится, то выигрывает Боря. Кто выигрывает при правильной игре?

\textit{Решение:} Заметим, что последним ходом Боря может поставить одно из четырёх чисел 1, 3, 7 или 9. Действительно, иначе полученное число сразу же будет делиться либо на 2 либо на 5. Давайте приведём выигрышную стратегию за Аню. 

На первый ход она ставит число 3. Если Боря своим ходом ставит одно из трёх чисел 1, 7 или 9, то Аня своими ходами просто "добирает"\  оставшиеся два числа и у Бори просто напросто не остается последнего хода. В ином случае Аня всегда ставит 9 на свой второй ход. Пусть Боря своими двумя ходами поставил числа $x$ и $y$. Достаточно показать, что какие бы два числа не взял Боря, Аня всегда может третьим ходом поставить такое число, что сумма получившихся пяти цифр будет давать остаток 2 при делении на 3 (так как последним ходом Боря ставит 1 или 7 и делает итоговое число кратным трём). Посмотрим на то, какие числа есть у Бори для первых двух ходов: $0, 1, 2, 4, 5, 6, 7, 8$. А если смотреть на остатки при делении на 3, то это $0, 1, 2, 1, 2, 0, 1, 2$. Теперь ясно видно, что если Ане понадобится число с остатком $r$ для дополнения хода Бори, то она такое обязательно найдёт $-$ все три единицы и двойки за два хода Боря не заберёт, а если а если заберёт оба нуля, то они нам и не нужны, так как теперь нужна двойка. 

\textbf{Более интересная постановка задачи.}

\problem{2}
 Аня и Боря играют в игру. Они по очереди (начинает Аня) выписывают по одной \textbf{ненулевой} цифре, пока не получится шестизначное число. При этом все выписанные цифры различны. Аня выигрывает, если полученное шестизначное число делится хотя бы на одно из чисел: 2, 3 или 5. Если этого не случится, то выигрывает Боря. Кто выигрывает при правильной игре?

 \newpage

\hspace{3mm}


 \textit{Решение:} Как и в прошлой формулировке задачи заметим, что последним ходом Боря может поставить одно из четырёх чисел 1, 3, 7 или 9. Обозначим эти числа \textit{особенными.} 

 $\textit{Необходимое условие Ани.}$ Аня обязательно за 3 своих хода должна взять хотя бы 2 особенных числа, причём если их ровно 2, то они сравнимы по модулю 3.
 
 Действительно, если Аня возьмёт только одно особенное число, то останутся хотя бы два с разными остатками при делении на 3. Боря может выбрать из них одно, которое нарушит делимость на 3.

 Далее разобьём задачу на два случая:

 \noindent1) \textbf{Первый ход Ани не особенный.} Если Аня первым ходом ставит не особенное число, то Боря заранее понимает, что следующие её два хода будут особенными, причём одинакового остатка при делении на 3. Значит на уже на третьей цифре Боря будет понимать, какие цифры у Ани (по модулю 3) и будет знать какой остаток будет давать последняя цифра (так как две "одинаковые"\ забрала Аня, остались две другие с другим остатком), а значит будет знать, какой остаток будет давать сумма всех цифр кроме четвёртой, которая как раз зависит от него. Два из трёх возможных остатков будут рушить делимость на 3 и хотя бы одна цифра с нужным остатком будет не использована к 4 ходу.

 \noindent2) \textbf{Первый ход Ани особенный.} Здесь нужно рассмотреть два симметричных случая $-$ первый ход 1/7 или первый ход 3/9.

2.1) \textbf{Первый ход Ани 1 или 7.}  Не умаляя общности, пусть первым ходом она поставила 1, на делимость на 3 это никак не влияет. Тогда Боря своим ставит 4 (в конце решения мы поймём почему именно 4). Теперь опять возможны два варианта: второе число Ани особенное или не особенное. Если это не особенное число, то мы попадаем в ситуацию с первым случаем. Боря уже понимает, что третье число будет особенным причём точно семёркой, а значит его последнее число либо 3 либо 9. То есть знает сумму всех цифр кроме четвёртой и сможет поставить туда то, что ему надо. Давайте разберём случай, когда второе число особенное. Если это 3 или 9, то в очередной раз Боря понимает, как число Ани будет третьим $-$ семёрка. Иначе в конце у Бори останется два числа с разным остатком и опять он к третьей цифре уже знает всю нужную информацию. Пусть Аня поставила 7 своим вторым ходом. Тогда последний ход Бори будет либо 3 либо 9 и всё зависит от 4 и 5 цифры. 

\newpage

\hspace{3mm}

\noindentЗаметим, что цифр дающих остаток 1 больше не осталось и тогда, если Боря поставит на 4 место двойку, сумма четвертой и пятой цифры теперь может давать остаток либо 1 (2+2) либо 2, а все остальные цифры в сумме дают остаток $1+4+7+0 \equiv 0 \ (mod \ 3)$ и значит сумма цифр не делится на 3.

2.2) \textbf{Первый ход Ани 3 или 9.} На самом то деле этот случай почти ничем не отличается. Единственное отличие в том, что теперь вместо 4 вторым ходом нужно поставить 6. Но почему мы выбрали именно такие числа? Давайте внимательно посмотрим на набор остатков среди особенных чисел: $1(1), 3(0), 7(1), 9(0).$ А какие есть неособенные числа с остатками 0 или 1? Вот тут и вся разница с предыдущей формулировкой. Так как в этой задаче цифры обязатяльено ненулевые, то есть ровно одно неособенное число с остатком 1 и одно неособенное число с остатком 0. Именно их и ставит Боря на свой второй ход.

\textbf{Как новому Боре победить старую Аню?} Теперь у Бори есть контрстратегия против Ани. Пусть Аня ставит первые два хода 3 и 9. Тогда Боря ставит 6, а после 9 числа кратные трём заканчиваются. Тогда Боря ставит на 4 место двойку и сумма цифр может давать остаток либо $3+6+9+2+1+1 = 1 (mod \ 3)$ либо $3+6+9+2+2+1 = 2 (mod \ 3)$.